\section{Speicherung und Weitergabe von Informationen im historischen Kontext}
\label{sec:einleitung}

Das Speichern und Weitergeben von Informationen ist zentraler Bestandteil der
menschlichen Kultur und führte in der Menschheitsgeschichte zu bahnbrechenden
Erfindungen und laufenden Innovationen.

In der Steinzeit wurden Informationen an Höhlenwände gemalt, ab dem 3.
Jahrtausend vor Christus mittels Keilschrift (siehe \autoref{fig:keilschrift})
in Steintafeln gemeißelt und in der Antike auf Papyrus gezeichnet. Mitte des 15.
Jahrhunderts revolutionierte dann Gutenberg die Informationsweitergabe mit der
Erfindung des Buchdrucks, der die Produktion von Büchern in hohen Stückzahlen
und zu geringen Kosten ermöglichte. Auch heute noch wird diese Technik für die
Reproduktion von Texten verwendet.

Ende des 19. Jahrhunderts gelang es erstmals, Musik zu speichern. Hierfür wurden
zunächst Schallplatten aus Hartgummi eingesetzt. Dieser wurde um 1900 dann durch
eine Pressmasse abgelöst, die im Wesentlichen aus Schellack\footnote{harzige
Ausscheidung von Lackschildläusen \cite{schellack}} (siehe
\autoref{fig:schallplatte-schellack}) bestand. Auf 12 Zoll-Schellackplatten
konnten Musikstücke mit ca. 4 Minuten Spielzeit pro Seite gespeichert werden.
Mit dem Einsatz des Kunststoffes Polyvinylchlorid ab den 50er Jahren des 20.
Jahrhunderts verbesserte sich die Tonqualität von Schallplatten deutlich.
Außerdem konnten längere Stücke gespeichert und wiedergegeben werden.
\cite{schallplatte1}

Mit der Compact Cassette (CC, siehe \autoref{fig:compact-cassette}), die 1963
auf den Markt kam, war es dann jedermann möglich, selbst Musik aufzunehmen und
dauerhaft zu speichern \cite{kassette1}. Compact Cassetten verwenden ein
Magnetband, das aus einer langen schmalen Kunststofffolie besteht, die mit einem
magnetisierbaren Material beschichtet wurde \cite{kassette2}.

In den 90er Jahren verdrängte  die Compact Disc (CD) innerhalb weniger Jahren
sowohl die Schallplatte als auch die Compact Cassette. Die CD zeichnet sich
durch ihre hohe Speicherkapazität sowie geringere Abnutzungserscheinungen aus
und wurde deshalb zum universellen Speichermedium für Musik, Dokumente, Bilder
und Filme. \cite{cd_durchbruch}

\begin{figure}[h]
    \begin{center}
        \begin{minipage}[t]{0.3\textwidth}
            \begin{center}
                \includegraphics[height=0.1\textheight]{Bilder/Einleitung/keilschrift.png}
                \caption[Altpersische Keilschrift \newline \url{http://www.faz.net/aktuell/feuilleton/geisteswissenschaften/keilschriften-zehntausend-freunde-mesopotamiens-1657070.html} (zuletzt aufgerufen am 19.09.2015)]{Altpersische Keilschrift}
                \label{fig:keilschrift}
            \end{center}
        \end{minipage}
        \hspace{0.025\textwidth}
        \begin{minipage}[t]{0.3\textwidth}
            \begin{center}
                \includegraphics[height=0.1\textheight]{Bilder/Einleitung/schallplatte-schellack.png}
                \caption[Schellackschallplatte \newline \url{https://de.wikipedia.org/wiki/Datei:Schallplatte_Deutsche_Grammophon_Stimme_seines_Herrn.jpg} (zuletzt aufgerufen am 19.09.2015)]{Schellackschall-platte}
                \label{fig:schallplatte-schellack}
            \end{center}
        \end{minipage}
        \hspace{0.025\textwidth}
        \begin{minipage}[t]{0.3\textwidth}
            \begin{center}
                \includegraphics[height=0.1\textheight]{Bilder/Einleitung/compact-cassette.png}
                \caption[Compact Cassette \newline \url{https://de.wikipedia.org/wiki/Datei:Compactcassette.jpg} (zuletzt aufgerufen am 19.09.2015)]{Compact Cassette}
                \label{fig:compact-cassette}
            \end{center}
        \end{minipage}
    \end{center}
\end{figure}

Kunststoffe sind - wie die oben aufgeführten Beispielen zeigen - für die
Informationstechnologie schon jetzt unersetzbar und kommen - neben der
Speicherung - auch bei der Übertragung großer Datenmengen zum Einsatz.

Sowohl für die Übertragung als auch für das Auslesen von Informationen werden
Lichtsignale verwendet. Hierfür werden vermehrt transparente Materialien
benötigt, die sich für den Alltagsgebrauch eignen.

Diese Arbeit behandelt den Einsatz von transparenten Kunststoffen in der
Informationstechnologie sowie ihre spezielle Eignung für die Speicherung und
Weitergabe von Informationen. Dies wird an den Beispielen der Compact Disc
(siehe \autoref{fig:compact-disc}) und der polymer optischen Faser (POF,  siehe
\autoref{fig:pof}) aufgezeigt. Dabei wird auf die Funktionsweise und Produktion
von optischen Datenträgern und optischen Wellenleitern eingangen. Außerdem
werden die physikalischen Eigenschaften und die Herstellung der verwendeten
Kunststoffe erläutert.

\begin{figure}[h]
    \begin{center}
        \begin{minipage}[t]{0.3\textwidth}
            \begin{center}
                \includegraphics[height=0.1\textheight]{Bilder/Einleitung/compact-disk.png}
                \caption[Compact Disc \newline \url{https://en.wikipedia.org/wiki/File:Compact_disc.svg} (zuletzt aufgerufen am 19.09.2015)]{Compact Disc}
                \label{fig:compact-disc}
            \end{center}
        \end{minipage}
        \hspace{0.025\textwidth}
        \begin{minipage}[t]{0.3\textwidth}
            \begin{center}
                \includegraphics[height=0.1\textheight]{Bilder/Einleitung/pof.png}
                \caption[polymer optische Faser \newline \url{http://www.heise.de/tr/imgs/08/2/5/4/2/1/5/c337ef89957e0f2b.jpg} (zuletzt aufgerufen am 19.09.2015)]{polymer optische Faser}
                \label{fig:pof}
            \end{center}
        \end{minipage}
    \end{center}
\end{figure}
